\structuralpart{ЗАКЛЮЧЕНИЕ}

Цель работы достигнута: рекурсивный и нерекурсивный алгоритмы определения
количества $k$-значных натуральных чисел с заданной суммой цифр сравнены
на основе графовых моделей.

В ходе выполнения работы решены следующие задачи:
\begin{itemize}
    \item даны определения операционного и информационного отношений, графа
    управления, информационного графа, операционной и информационной историй;
    \item разработаны нерекурсивный алгоритм заполнения таблицы значений
    и рекурсивный алгоритм без циклов, массивов, строк и срезов;
    \item алгоритмы реализованы на языке Свифт (Swift);
    \item реализации прошли все 14 функциональных проверок, по 7 для каждой
    реализации;
    \item каждый алгоритм описан четырьмя графовыми моделями; истории построены
    для $k=1$, $s=1$;
    \item показано, что в нерекурсивном алгоритме параллельно могут вычисляться
    ячейки одной строки таблицы, а строки~--- только последовательно;
    в рекурсивном алгоритме независимы вызовы $F(i-1,j-c)$ и $S(i,j,c+1)$,
    а сумма их результатов выполняется после обоих вызовов.
\end{itemize}

По графовым моделям рекурсивный алгоритм требует значительно большего числа
срабатываний операторов из-за повторного вычисления одинаковых значений,
а число ярусов информационных историй обоих алгоритмов близко.
